%%%%%%%%%%%%%%%%%%%%%%%%%%%%%%%%%%%%%%%%%%%%%%%%%%%%%%%%%%%%%%%%%%%%%%%%%%%%%%%%%%
\begin{frame}[fragile]\frametitle{}
\begin{center}
{\Large Case Study: Titanic Disaster}

% % (Ref: https://github.com/donnemartin/data-science-ipython-notebooks/blob/master/kaggle/titanic.ipynb)
\end{center}
\end{frame}

%%%%%%%%%%%%%%%%%%%%%%%%%%%%%%%%%%%%%%%%%%%%%%%%%%%
\begin{frame}[fragile]\frametitle{The Disaster}
\begin{center}
\includegraphics[width=0.85\linewidth,keepaspectratio]{titanic}
\end{center}
\tiny{(Image Reference: https://en.wikipedia.org/wiki/Willy\_Stower)}
\end{frame}

%%%%%%%%%%%%%%%%%%%%%%%%%%%%%%%%%%%%%%%%%%%%%%%%%%%
\begin{frame}[fragile]\frametitle{The Titanic Disaster: Background}
\begin{itemize}
\item The sinking of the Titanic is one of the most infamous shipwrecks in history.
\item It sank after colliding with an iceberg, killing 1502 out of 2224 passengers and crew.
\item One of the reasons for loss of life was that there were not enough lifeboats for the passengers and crew. 
\item Some groups of people were more likely to survive than others, such as women, children, and the upper-class.
\end{itemize}
\end{frame}

%%%%%%%%%%%%%%%%%%%%%%%%%%%%%%%%%%%%%%%%%%%%%%%%%%%%%%%%%%%%%%%%%%%%%%%%%%%%%%%%%%
% Folded away (Sep 2026 review): section divider; the Kaggle context is now two frames.
% \begin{frame}[fragile]\frametitle{}
% \begin{center}
% {\Large Kaggle}
% \end{center}
% \end{frame}


%%%%%%%%%%%%%%%%%%%%%%%%%%%%%%%%%%%%%%%%%%%%%%%%%%%
% Folded away (Sep 2026 review): screenshot of the competition page.
% \begin{frame}[fragile]\frametitle{Kaggle Competition}
% \begin{center}
% 
\includegraphics[width=0.65\linewidth,keepaspectratio]{titanic_kaggle}
% \end{center}
% \end{frame}

%%%%%%%%%%%%%%%%%%%%%%%%%%%%%%%%%%%%%%%%%%%%%%%%%%%
\begin{frame}[fragile]\frametitle{What is Kaggle?}
\begin{itemize}
\item Kaggle is a site where people create algorithms and compete against machine learning practitioners around the world. 
\item Your algorithm wins the competition if it's the most accurate on a particular data set. 
\item Kaggle is a fun way to practice your machine learning skills.
\item Indian counterpart is ``Analytics Vidhya''. 
\item Many competitions and good prizes.
\item The Titanic competition is a beginner favourite: use passenger details to predict who survived.
\end{itemize}
\end{frame}


%%%%%%%%%%%%%%%%%%%%%%%%%%%%%%%%%%%%%%%%%%%%%%%%%%%
% Folded into 'What is Kaggle?' (Sep 2026 review).
% \begin{frame}[fragile]\frametitle{The Titanic Competition: Goal}
% \begin{itemize}
% \item Kaggle has created a number of competitions designed for beginners.
% \item One of the most popular of these competitions is Titanic competition
% \item In this competition, we have a data set of different information about passengers onboard the Titanic, and we see if we can use that information to predict whether those people survived or not.
% \end{itemize}
% \end{frame}


%%%%%%%%%%%%%%%%%%%%%%%%%%%%%%%%%%%%%%%%%%%%%%%%%%%
\begin{frame}[fragile]\frametitle{The Titanic Competition: Train vs Test Data}
\begin{itemize}
\item Each Kaggle competition has two key data files that you will work with - a training set and a testing set.
\item The training set contains data we can use to train our model. 
\item It has a number of feature columns which contain various descriptive data, as well as a column of the target values we are trying to predict: in this case, Survival.
\item The testing set contains all of the same feature columns, but is missing the target value column. Additionally, the testing set usually has fewer observations (rows) than the training set.
\end{itemize}
\end{frame}

%%%%%%%%%%%%%%%%%%%%%%%%%%%%%%%%%%%%%%%%%%%%%%%%%%%%%%%%%%%%%%%%%%%%%%%%%%%%%%%%%%
\begin{frame}[fragile]\frametitle{}
\begin{center}
{\Large Data}
\end{center}
\end{frame}


%%%%%%%%%%%%%%%%%%%%%%%%%%%%%%%%%%%%%%%%%%%%%%%%%%%
\begin{frame}[fragile]\frametitle{Data Dictionary}
Below are the descriptions contained in that data dictionary
\begin{itemize}
\item PassengerID: A column added by Kaggle to identify each row
\item Survived: Passenger survived or not (0=No, 1=Yes)
\item Pclass: The class of the ticket (1=1st, 2=2nd, 3=3rd)
\item Sex: The passenger's sex
\item Age: The passenger's age in years
\item SibSp: The number of siblings or spouses 
\item Parch: The number of parents or children
\item Ticket: The passenger's ticket number
\item Fare: The fare the passenger paid
\item Cabin: The passenger's cabin number
\item Embarked: The port where the passenger embarked (C=Cherbourg, Q=Queenstown, S=Southampton)
\end{itemize}
\end{frame}



%%%%%%%%%%%%%%%%%%%%%%%%%%%%%%%%%%%%%%%%%%%%%%%%%%%
\begin{frame}[fragile]\frametitle{Load the data}
Lines starting with \texttt{>>>} are typed; everything below is what Python prints back.
\begin{lstlisting}
>>> import pandas as pd
>>> df_test = pd.read_csv("titanic_test.csv")
>>> df_train = pd.read_csv("titanic_train.csv")
>>> print("Dimensions of train: {}".format(df_train.shape))
Dimensions of train: (891, 12)
>>> print("Dimensions of test: {}".format(df_test.shape))
Dimensions of test: (418, 11)
\end{lstlisting}
\end{frame}

%%%%%%%%%%%%%%%%%%%%%%%%%%%%%%%%%%%%%%%%%%%%%%%%%%%
\begin{frame}[fragile]\frametitle{Setup: Imports and Plot Settings}
The plotting frames below use these names:
\begin{lstlisting}
import numpy as np
import matplotlib.pyplot as plt

figsize_with_subplots = (10, 10)   # figure size for grids of plots
bin_size = 10                      # width of an age bin, in years
\end{lstlisting}
\end{frame}

%%%%%%%%%%%%%%%%%%%%%%%%%%%%%%%%%%%%%%%%%%%%%%%%%%%%%%%%%%%%%%%%%%%%%%%%%%%%%%%%%%
\begin{frame}[fragile]\frametitle{}
\begin{center}
{\Large Data Exploration}
\end{center}
\end{frame}

%%%%%%%%%%%%%%%%%%%%%%%%%%%%%%%%%%%%%%%%%%%%%%%%%%%
\begin{frame}[fragile]\frametitle{head()}
Let's take a look at the first few rows of the train dataframe.

\begin{center}
\includegraphics[width=\linewidth,keepaspectratio]{titk1}
\end{center}

\begin{lstlisting}
df_train.head()
\end{lstlisting}

\end{frame}

%%%%%%%%%%%%%%%%%%%%%%%%%%%%%%%%%%%%%%%%%%%%%%%%%%%
% Optional (Sep 2026 review): tail() shows the same columns as head().
% \begin{frame}[fragile]\frametitle{Exploring the Data: tail()}
% Let's take a look at the last few rows of the train dataframe.

% \begin{center}
% \includegraphics[width=\linewidth,keepaspectratio]{titk2}
% \end{center}

% \begin{lstlisting}
% df_train.tail()
% \end{lstlisting}
% \end{frame}

%%%%%%%%%%%%%%%%%%%%%%%%%%%%%%%%%%%%%%%%%%%%%%%%%%%
\begin{frame}[fragile]\frametitle{Column Data Types}
\begin{itemize}
\item Type 'object' is a string for pandas, which poses problems with machine learning algorithms.
\item If we want to use these as features, we'll need to convert these to number representations.
\end{itemize}

Lines starting with \texttt{>>>} are typed; everything below is what Python prints back.
\begin{lstlisting}
>>> df_train.dtypes
PassengerId      int64
Survived         int64
Pclass           int64
Name            object
Sex             object
Age            float64
SibSp            int64
Parch            int64
Ticket          object
Fare           float64
Cabin           object
Embarked        object
\end{lstlisting}

\end{frame}

%%%%%%%%%%%%%%%%%%%%%%%%%%%%%%%%%%%%%%%%%%%%%%%%%%%
\begin{frame}[fragile]\frametitle{Dataframe information}
\begin{itemize}
\item Age, Cabin, and Embarked are missing values.
\item Cabin has too many missing values.
\end{itemize}

\begin{lstlisting}
>>> df_train.info()
Data columns (total 12 columns):
PassengerId    891 non-null int64
Survived       891 non-null int64
Pclass         891 non-null int64
Name           891 non-null object
Sex            891 non-null object
Age            714 non-null float64
SibSp          891 non-null int64
Parch          891 non-null int64
Ticket         891 non-null object
Fare           891 non-null float64
Cabin          204 non-null object
Embarked       889 non-null object
\end{lstlisting}

\end{frame}


%%%%%%%%%%%%%%%%%%%%%%%%%%%%%%%%%%%%%%%%%%%%%%%%%%%
\begin{frame}[fragile]\frametitle{Descriptive Statistics}
Run \texttt{df\_train.describe()}:
\begin{itemize}
\item The mean of Survived is 0.38: about 38\% of the passengers survived
\item Age: 714 of 891 rows are present; mean 29.7, median 28
\item Fare: mean 32.2 but median 14.45, so a few very high fares stretch the right tail
\end{itemize}
\begin{center}
\includegraphics[width=\linewidth,keepaspectratio]{titk3}
\end{center}
\end{frame}

%%%%%%%%%%%%%%%%%%%%%%%%%%%%%%%%%%%%%%%%%%%%%%%%%%%
\begin{frame}[fragile]\frametitle{Survival Counts}
Plot death and survival counts

\begin{center}
\includegraphics[width=0.4\linewidth,keepaspectratio]{titk4}
\end{center}

\begin{lstlisting}
# Set up a grid of plots
fig = plt.figure(figsize=figsize_with_subplots) 
fig_dims = (3, 2)

plt.subplot2grid(fig_dims, (0, 0))
df_train['Survived'].value_counts().plot(kind='bar', 
                                         title='Death and Survival Counts')
\end{lstlisting}

\end{frame}

%%%%%%%%%%%%%%%%%%%%%%%%%%%%%%%%%%%%%%%%%%%%%%%%%%%
\begin{frame}[fragile]\frametitle{Class Counts}
Plot Pclass counts

\begin{center}
\includegraphics[width=0.5\linewidth,keepaspectratio]{titk5}
\end{center}

\begin{lstlisting}
plt.subplot2grid(fig_dims, (0, 1))
df_train['Pclass'].value_counts().plot(kind='bar', 
                                       title='Passenger Class Counts')
\end{lstlisting}

\end{frame}

%%%%%%%%%%%%%%%%%%%%%%%%%%%%%%%%%%%%%%%%%%%%%%%%%%%
\begin{frame}[fragile]\frametitle{Sex, Embarked, Age}
Plot Sex counts
\begin{lstlisting}
plt.subplot2grid(fig_dims, (1, 0))
df_train['Sex'].value_counts().plot(kind='bar', 
                                    title='Gender Counts')
plt.xticks(rotation=0)

# Plot Embarked counts

plt.subplot2grid(fig_dims, (1, 1))
df_train['Embarked'].value_counts().plot(kind='bar', 
                                         title='Ports of Embarkation Counts')

# Plot the Age histogram

plt.subplot2grid(fig_dims, (2, 0))
df_train['Age'].hist()
plt.title('Age Histogram')
\end{lstlisting}
\end{frame}

%%%%%%%%%%%%%%%%%%%%%%%%%%%%%%%%%%%%%%%%%%%%%%%%%%%
\begin{frame}[fragile]\frametitle{Cross Tab}
We'll determine which proportion of passengers survived based on their passenger class. Generate a cross tab of Pclass and Survived:

\begin{center}
\includegraphics[width=0.3\linewidth,keepaspectratio]{titk6}
\end{center}


\begin{lstlisting}
pclass_xt = pd.crosstab(df_train['Pclass'], df_train['Survived'])
pclass_xt
\end{lstlisting}

\end{frame}


%%%%%%%%%%%%%%%%%%%%%%%%%%%%%%%%%%%%%%%%%%%%%%%%%%%
\begin{frame}[fragile]\frametitle{Map Sex to Numbers}
Generate a mapping of Sex from a string to a number representation:
\begin{lstlisting}
>>> sexes = sorted(df_train['Sex'].unique())
>>> genders_mapping = dict(zip(sexes, range(len(sexes))))
>>> genders_mapping
{'female': 0, 'male': 1}
\end{lstlisting}
\end{frame}


%%%%%%%%%%%%%%%%%%%%%%%%%%%%%%%%%%%%%%%%%%%%%%%%%%%
\begin{frame}[fragile]\frametitle{Apply the Sex Mapping}
Transform Sex from a string to a number representation:

\begin{center}
\includegraphics[width=0.8\linewidth,keepaspectratio]{titk7}
\end{center}


\begin{lstlisting}
df_train['Sex_Val'] = df_train['Sex'].map(genders_mapping).astype(int)
df_train.head()
\end{lstlisting}

\end{frame}

%%%%%%%%%%%%%%%%%%%%%%%%%%%%%%%%%%%%%%%%%%%%%%%%%%%
\begin{frame}[fragile]\frametitle{Survival by Sex}
Plot a normalized cross tab for Sex\_Val and Survived:
\begin{center}
\includegraphics[width=0.4\linewidth,keepaspectratio]{titk8}
\end{center}

The majority of females survived (74\%), but not males (19\%).


\begin{lstlisting}
sex_val_xt = pd.crosstab(df_train['Sex_Val'], df_train['Survived'])
sex_val_xt_pct = sex_val_xt.div(sex_val_xt.sum(1).astype(float), axis=0)
sex_val_xt_pct.plot(kind='bar', stacked=True, title='Survival Rate by Gender')
\end{lstlisting}

\end{frame}

%%%%%%%%%%%%%%%%%%%%%%%%%%%%%%%%%%%%%%%%%%%%%%%%%%%
\begin{frame}[fragile]\frametitle{Sex Counts by Class}
Any insights by looking at both Sex and Pclass? Count males and females in each Pclass:
\begin{lstlisting}
>>> passenger_classes = sorted(df_train['Pclass'].unique())
>>> for p_class in passenger_classes:
...     print('M: ', p_class, len(df_train[(df_train['Sex'] == 'male') &
...                              (df_train['Pclass'] == p_class)]))
...     print('F: ', p_class, len(df_train[(df_train['Sex'] == 'female') &
...                              (df_train['Pclass'] == p_class)]))
...
M:  1 122
F:  1 94
M:  2 108
F:  2 76
M:  3 347
F:  3 144
\end{lstlisting}
\end{frame}

%%%%%%%%%%%%%%%%%%%%%%%%%%%%%%%%%%%%%%%%%%%%%%%%%%%
\begin{frame}[fragile]\frametitle{Female Survival by Class}
Plot survival rate by Sex:
\begin{lstlisting}
# Plot survival rate by Sex
females_df = df_train[df_train['Sex'] == 'female']
females_xt = pd.crosstab(females_df['Pclass'], df_train['Survived'])
females_xt_pct = females_xt.div(females_xt.sum(1).astype(float), axis=0)
females_xt_pct.plot(kind='bar',  stacked=True, 
                    title='Female Survival Rate by Passenger Class')
plt.xlabel('Passenger Class')
plt.ylabel('Survival Rate')
\end{lstlisting}

\end{frame}

%%%%%%%%%%%%%%%%%%%%%%%%%%%%%%%%%%%%%%%%%%%%%%%%%%%
\begin{frame}[fragile]\frametitle{Male Survival by Class}
Plot survival rate by Pclass:

The vast majority of females in First and Second class survived. Males in First class had the highest chance for survival.

\begin{lstlisting}
# Plot survival rate by Pclass
males_df = df_train[df_train['Sex'] == 'male']
males_xt = pd.crosstab(males_df['Pclass'], df_train['Survived'])
males_xt_pct = males_xt.div(males_xt.sum(1).astype(float), axis=0)
males_xt_pct.plot(kind='bar',    stacked=True, 
                  title='Male Survival Rate by Passenger Class')
plt.xlabel('Passenger Class')
plt.ylabel('Survival Rate')                      
\end{lstlisting}

\end{frame}


%%%%%%%%%%%%%%%%%%%%%%%%%%%%%%%%%%%%%%%%%%%%%%%%%%%
% Optional (Sep 2026 review): the Embarked chain is dropped again at Final Data Prep, so it is skipped in the session. Restore all six frames together.
% \begin{frame}[fragile]\frametitle{Cleaning the Data: Missing Embarked Rows}
% The Embarked column might be an important feature but it is missing a couple data points which might pose a problem for machine learning algorithms:

% \begin{center}
% \includegraphics[width=0.8\linewidth,keepaspectratio]{titk9}
% \end{center}

% \begin{lstlisting}
% df_train[df_train['Embarked'].isnull()]
% \end{lstlisting}

% \end{frame}


%%%%%%%%%%%%%%%%%%%%%%%%%%%%%%%%%%%%%%%%%%%%%%%%%%%
% Optional (Sep 2026 review): the Embarked chain is dropped again at Final Data Prep, so it is skipped in the session. Restore all six frames together.
% \begin{frame}[fragile]\frametitle{Cleaning the Data: Fill Missing Embarked}
% Since the vast majority of passengers embarked at 'S' (Southampton), we assign the 2 missing values in Embarked to 'S':
% \begin{lstlisting}
% >>> df_train['Embarked'] = df_train['Embarked'].fillna('S')
% >>> df_train['Embarked'].isnull().sum()
% 0
% \end{lstlisting}
% \end{frame}

%%%%%%%%%%%%%%%%%%%%%%%%%%%%%%%%%%%%%%%%%%%%%%%%%%%
% Optional (Sep 2026 review): the Embarked chain is dropped again at Final Data Prep, so it is skipped in the session. Restore all six frames together.
% \begin{frame}[fragile]\frametitle{Cleaning the Data: Map Embarked to Numbers}
% Prepare to map Embarked from a string to a number representation:
% \begin{lstlisting}
% >>> embarked_locs = sorted(df_train['Embarked'].unique())
% >>> embarked_locs_mapping = dict(zip(embarked_locs,
% ...                                  range(1, len(embarked_locs) + 1)))
% >>> embarked_locs_mapping
% {'C': 1, 'Q': 2, 'S': 3}
% \end{lstlisting}
% \end{frame}

%%%%%%%%%%%%%%%%%%%%%%%%%%%%%%%%%%%%%%%%%%%%%%%%%%%
% Optional (Sep 2026 review): the Embarked chain is dropped again at Final Data Prep, so it is skipped in the session. Restore all six frames together.
% \begin{frame}[fragile]\frametitle{Cleaning the Data: Apply the Embarked Mapping}
% Transform Embarked from a string to a number representation to prepare it for machine learning algorithms:

% \begin{center}
% \includegraphics[width=0.6\linewidth,keepaspectratio]{titk10}
% \end{center}


% \begin{lstlisting}
% df_train['Embarked_Val'] = df_train['Embarked'] \
%                                .map(embarked_locs_mapping) \
%                                .astype(int)
% df_train.head()
% \end{lstlisting}

% \end{frame}

%%%%%%%%%%%%%%%%%%%%%%%%%%%%%%%%%%%%%%%%%%%%%%%%%%%
% Superseded (Sep 2026): the old order mapped NaN to 0 (sorted() fails on NaN in Python 3) then replaced it; merged into the frames above.
% \begin{frame}[fragile]\frametitle{Cleaning the Data: Fill Missing Embarked}
% Since the vast majority of passengers embarked in 'S': 3, we assign the missing values in Embarked to 'S':
% \begin{lstlisting}
% if len(df_train[df_train['Embarked'].isnull()]) > 0:
%     df_train.replace({'Embarked_Val' : 
%                    { embarked_locs_mapping[nan] : embarked_locs_mapping['S'] 
%                    }
%                }, 
%                inplace=True)
% \end{lstlisting}
% \end{frame}

%%%%%%%%%%%%%%%%%%%%%%%%%%%%%%%%%%%%%%%%%%%%%%%%%%%
% Superseded: merged into "Fill Missing Embarked".
% \begin{frame}[fragile]\frametitle{Cleaning the Data: Verify No Missing Values}
% Verify we do not have any more NaNs for Embarked\_Val:
% \begin{lstlisting}
% embarked_locs = sorted(df_train['Embarked_Val'].unique())
% embarked_locs
% 
% array([1, 2, 3])
% \end{lstlisting}
% \end{frame}
% Old histogram frame (image titk11 showed the NaN bin as 0), superseded:
% \begin{frame}[fragile]\frametitle{Cleaning the Data: Embarked Histogram}
% Plot the histogram for Embarked\_Val:
% 
% \begin{center}
% \includegraphics[width=0.6\linewidth,keepaspectratio]{titk11}
% \end{center}
% 
% \begin{lstlisting}
% df_train['Embarked_Val'].hist(bins=len(embarked_locs), range=(0, 3))
% plt.title('Port of Embarkation Histogram')
% plt.xlabel('Port of Embarkation')
% plt.ylabel('Count')
% plt.show()
% \end{lstlisting}
% 
% 
% \end{frame}

%%%%%%%%%%%%%%%%%%%%%%%%%%%%%%%%%%%%%%%%%%%%%%%%%%%
% Optional (Sep 2026 review): the Embarked chain is dropped again at Final Data Prep, so it is skipped in the session. Restore all six frames together.
% \begin{frame}[fragile]\frametitle{Exploring the Data: Survival by Embarkation Setup}
% Plot a normalized cross tab for Embarked\_Val and Survived:
% \begin{lstlisting}
% embarked_val_xt = pd.crosstab(df_train['Embarked_Val'], df_train['Survived'])
% embarked_val_xt_pct = \
%     embarked_val_xt.div(embarked_val_xt.sum(1).astype(float), axis=0)
% embarked_val_xt_pct.plot(kind='bar', stacked=True)
% plt.title('Survival Rate by Port of Embarkation')
% plt.xlabel('Port of Embarkation')
% plt.ylabel('Survival Rate')
% \end{lstlisting}
% \end{frame}

%%%%%%%%%%%%%%%%%%%%%%%%%%%%%%%%%%%%%%%%%%%%%%%%%%%
% Optional (Sep 2026 review): the Embarked chain is dropped again at Final Data Prep, so it is skipped in the session. Restore all six frames together.
% \begin{frame}[fragile]\frametitle{Exploring the Data: Survival by Embarkation Result}
% \begin{center}
% \includegraphics[width=0.6\linewidth,keepaspectratio]{titk12}
% \end{center}
% It appears those that embarked in location 'C': 1 had the highest rate of survival. 
% \end{frame}

%%%%%%%%%%%%%%%%%%%%%%%%%%%%%%%%%%%%%%%%%%%%%%%%%%%
\begin{frame}[fragile]\frametitle{}
\begin{center}
{\Large Feature Engineering: Age and Family Size}
\end{center}
\end{frame}

%%%%%%%%%%%%%%%%%%%%%%%%%%%%%%%%%%%%%%%%%%%%%%%%%%%
\begin{frame}[fragile]\frametitle{Missing Age Values}
The Age column seems like an important feature: unfortunately it is missing many values. Filter to view missing Age values:

\begin{center}
\includegraphics[width=0.25\linewidth,keepaspectratio]{titk13}
\end{center}


\begin{lstlisting}
df_train[df_train['Age'].isnull()][['Sex', 'Pclass', 'Age']].head()
\end{lstlisting}

\end{frame}

%%%%%%%%%%%%%%%%%%%%%%%%%%%%%%%%%%%%%%%%%%%%%%%%%%%
\begin{frame}[fragile]\frametitle{Fill Missing Ages}
Determine the Age typical for each passenger class by Sex\_Val. We'll use the median instead of the mean because the Age histogram seems to be right skewed.
\begin{lstlisting}
# Keep Age intact: AgeFill is a copy where each missing age is replaced
# by the median age of that (Sex_Val, Pclass) group
df_train['AgeFill'] = df_train.groupby(['Sex_Val', 'Pclass'])['Age'] \
                              .transform(lambda x: x.fillna(x.median()))
\end{lstlisting}
\end{frame}

%%%%%%%%%%%%%%%%%%%%%%%%%%%%%%%%%%%%%%%%%%%%%%%%%%%
\begin{frame}[fragile]\frametitle{Age Histogram}
Plot a histogram of AgeFill, split by Survived:
\begin{lstlisting}
# Set up a grid of plots
fig, axes = plt.subplots(2, 1, figsize=figsize_with_subplots)

# Histogram of AgeFill segmented by Survived
df1 = df_train[df_train['Survived'] == 0]['AgeFill']
df2 = df_train[df_train['Survived'] == 1]['AgeFill']
max_age = int(max(df_train['AgeFill']))
axes[0].hist([df1, df2],
             bins=max_age // bin_size,
             range=(1, max_age),
             stacked=True)
axes[0].legend(('Died', 'Survived'), loc='best')
axes[0].set_title('Survivors by Age Groups Histogram')
axes[0].set_xlabel('Age')
axes[0].set_ylabel('Count')
\end{lstlisting}
\end{frame}

%%%%%%%%%%%%%%%%%%%%%%%%%%%%%%%%%%%%%%%%%%%%%%%%%%%
% Kept as its own frame (Sep 2026 review): merging it with the 15-line code above overflows the slide.
\begin{frame}[fragile]\frametitle{Age Histogram Result}
\begin{center}
\includegraphics[width=0.6\linewidth,keepaspectratio]{titk14_agefill}
\end{center}
\end{frame}

%%%%%%%%%%%%%%%%%%%%%%%%%%%%%%%%%%%%%%%%%%%%%%%%%%%
\begin{frame}[fragile]\frametitle{Age vs Survival Scatter}
Plot AgeFill against Survived:
\begin{center}
\includegraphics[width=0.5\linewidth,keepaspectratio]{titk15_agefill}
\end{center}

Unfortunately, the graphs above do not seem to clearly show any insights.


\begin{lstlisting}
axes[1].scatter(df_train['Survived'], df_train['AgeFill'])
axes[1].set_title('Survivors by Age Plot')
axes[1].set_xlabel('Survived')
axes[1].set_ylabel('Age')
\end{lstlisting}

\end{frame}

%%%%%%%%%%%%%%%%%%%%%%%%%%%%%%%%%%%%%%%%%%%%%%%%%%%
\begin{frame}[fragile]\frametitle{Age Density by Class}
Plot AgeFill density by Pclass:

\begin{center}
\includegraphics[width=0.6\linewidth,keepaspectratio]{titk16}
\end{center}


\begin{lstlisting}
for pclass in passenger_classes:
    df_train.AgeFill[df_train.Pclass == pclass].plot(kind='kde')
plt.title('Age Density Plot by Passenger Class')
plt.xlabel('Age')
plt.legend(('1st Class', '2nd Class', '3rd Class'), loc='best')
\end{lstlisting}

\end{frame}


%%%%%%%%%%%%%%%%%%%%%%%%%%%%%%%%%%%%%%%%%%%%%%%%%%%
\begin{frame}[fragile]\frametitle{Age and Class Insights}
\begin{itemize}
\item The first class passengers were generally older then second class passengers, which in turn were older than third class passengers. 
\item We've determined that first class passengers had a higher survival rate than second class passengers, which in turn had a higher survival rate than third class passengers.
\end{itemize}
\end{frame}

%%%%%%%%%%%%%%%%%%%%%%%%%%%%%%%%%%%%%%%%%%%%%%%%%%%
\begin{frame}[fragile]\frametitle{Survivors' Age Histogram}
Most survivors come from the 20's to 30's age ranges, which the next two graphs help to explain.
\begin{center}
\includegraphics[width=0.8\linewidth,height=2.6cm,keepaspectratio]{titk17}
\end{center}
\begin{lstlisting}
# Set up a grid of plots
fig = plt.figure(figsize=figsize_with_subplots)
fig_dims = (3, 1)

# Plot the AgeFill histogram for Survivors
plt.subplot2grid(fig_dims, (0, 0))
survived_df = df_train[df_train['Survived'] == 1]
survived_df['AgeFill'].hist(bins=max_age // bin_size, range=(1, max_age))
\end{lstlisting}
\end{frame}

%%%%%%%%%%%%%%%%%%%%%%%%%%%%%%%%%%%%%%%%%%%%%%%%%%%
% Merged into 'Feature Engineering: Survivors' Age Histogram' (Sep 2026 review): picture now sits above the code.
% \begin{frame}[fragile]\frametitle{Feature Engineering: Survivors' Age Histogram Result}
% \begin{center}
% \includegraphics[width=0.8\linewidth,keepaspectratio]{titk17}
% \end{center}
% We see that most survivors come from the 20's to 30's age ranges and might be explained by the following two graphs. 
% \end{frame}

%%%%%%%%%%%%%%%%%%%%%%%%%%%%%%%%%%%%%%%%%%%%%%%%%%%
\begin{frame}[fragile]\frametitle{Surviving Females' Age Histogram}
Most surviving females are within their 20's.
\begin{center}
\includegraphics[width=0.8\linewidth,height=2.6cm,keepaspectratio]{titk18}
\end{center}
\begin{lstlisting}
# Plot the AgeFill histogram for Females
plt.subplot2grid(fig_dims, (1, 0))
females_df = df_train[(df_train['Sex_Val'] == 0) & (df_train['Survived'] == 1)]
females_df['AgeFill'].hist(bins=max_age // bin_size, range=(1, max_age))
\end{lstlisting}
\end{frame}

%%%%%%%%%%%%%%%%%%%%%%%%%%%%%%%%%%%%%%%%%%%%%%%%%%%
% Merged into 'Feature Engineering: Surviving Females' Age Histogram' (Sep 2026 review): picture now sits above the code.
% \begin{frame}[fragile]\frametitle{Feature Engineering: Females' Age Histogram Result}
% \begin{center}
% \includegraphics[width=0.8\linewidth,keepaspectratio]{titk18}
% \end{center}

% Shows most females are within their 20's.
% \end{frame}


%%%%%%%%%%%%%%%%%%%%%%%%%%%%%%%%%%%%%%%%%%%%%%%%%%%
\begin{frame}[fragile]\frametitle{Surviving First-Class Age Histogram}
Plots

\begin{center}
\includegraphics[width=0.8\linewidth,keepaspectratio]{titk19}
\end{center}

Shows most surviving first class passengers are within their 30's.

\begin{lstlisting}
# Plot the AgeFill histogram for first class passengers
plt.subplot2grid(fig_dims, (2, 0))
class1_df = df_train[(df_train['Pclass'] == 1) & (df_train['Survived'] == 1)]
class1_df['AgeFill'].hist(bins=max_age // bin_size, range=(1, max_age))
\end{lstlisting}

\end{frame}



%%%%%%%%%%%%%%%%%%%%%%%%%%%%%%%%%%%%%%%%%%%%%%%%%%%
\begin{frame}[fragile]\frametitle{Create FamilySize}
Define a new feature FamilySize that is the sum of Parch (number of parents or children on board) and SibSp (number of siblings or spouses):
\begin{center}
\includegraphics[width=0.9\linewidth,height=4.2cm,keepaspectratio]{titk20}
\end{center}
\begin{lstlisting}
df_train['FamilySize'] = df_train['SibSp'] + df_train['Parch']
df_train.head()
\end{lstlisting}
\end{frame}


%%%%%%%%%%%%%%%%%%%%%%%%%%%%%%%%%%%%%%%%%%%%%%%%%%%
% Merged into 'Feature Engineering: Create FamilySize' (Sep 2026 review): picture now sits above the code.
% \begin{frame}[fragile]\frametitle{Feature Engineering: FamilySize Column Added}
% \begin{center}
% \includegraphics[width=\linewidth,keepaspectratio]{titk20}
% \end{center}
% \end{frame}

%%%%%%%%%%%%%%%%%%%%%%%%%%%%%%%%%%%%%%%%%%%%%%%%%%%
\begin{frame}[fragile]\frametitle{FamilySize Histogram}
Plot a histogram of FamilySize:

\begin{center}
\includegraphics[width=0.8\linewidth,keepaspectratio]{titk21}
\end{center}


\begin{lstlisting}
df_train['FamilySize'].hist()
plt.title('Family Size Histogram')
\end{lstlisting}

\end{frame}

%%%%%%%%%%%%%%%%%%%%%%%%%%%%%%%%%%%%%%%%%%%%%%%%%%%
\begin{frame}[fragile]\frametitle{FamilySize vs Survival}
Plot a histogram of FamilySize segmented by Survived. It is not immediately obvious what impact FamilySize has on survival.
\begin{center}
\includegraphics[height=2.8cm,keepaspectratio]{titk22}
\end{center}
\begin{lstlisting}
family_sizes = sorted(df_train['FamilySize'].unique())
family_size_max = max(family_sizes)
df1 = df_train[df_train['Survived'] == 0]['FamilySize']
df2 = df_train[df_train['Survived'] == 1]['FamilySize']
plt.hist([df1, df2], 
         bins=family_size_max + 1, 
         range=(0, family_size_max), 
         stacked=True)
plt.legend(('Died', 'Survived'), loc='best')
plt.title('Survivors by Family Size')
\end{lstlisting}
\end{frame}

%%%%%%%%%%%%%%%%%%%%%%%%%%%%%%%%%%%%%%%%%%%%%%%%%%%
% Merged into 'Feature Engineering: FamilySize vs Survival' (Sep 2026 review): picture now sits above the code.
% \begin{frame}[fragile]\frametitle{Feature Engineering: FamilySize vs Survival Result}
% \begin{center}
% \includegraphics[width=0.8\linewidth,keepaspectratio]{titk22}
% \end{center}
% It is not immediately obvious what impact FamilySize has on survival.
% \end{frame}

%%%%%%%%%%%%%%%%%%%%%%%%%%%%%%%%%%%%%%%%%%%%%%%%%%%
\begin{frame}[fragile]\frametitle{}
\begin{center}
{\Large Final Data Prep}
\end{center}
\end{frame}

%%%%%%%%%%%%%%%%%%%%%%%%%%%%%%%%%%%%%%%%%%%%%%%%%%%
\begin{frame}[fragile]\frametitle{Find String Columns}
Many machine learning algorithms do not work on strings and they usually require the data to be in an array, not a DataFrame. Show only the columns of type 'object' (strings):

\begin{lstlisting}
>>> df_train.dtypes[df_train.dtypes.map(lambda x: x == 'object')]
Name        object
Sex         object
Ticket      object
Cabin       object
Embarked    object
dtype: object
\end{lstlisting}
\end{frame}

%%%%%%%%%%%%%%%%%%%%%%%%%%%%%%%%%%%%%%%%%%%%%%%%%%%
\begin{frame}[fragile]\frametitle{Drop String Columns}
Drop the columns we won't use:
\begin{lstlisting}
df_train = df_train.drop(['Name', 'Sex', 'Ticket', 'Cabin', 'Embarked'],
                         axis=1)
\end{lstlisting}

\end{frame}

%%%%%%%%%%%%%%%%%%%%%%%%%%%%%%%%%%%%%%%%%%%%%%%%%%%
\begin{frame}[fragile]\frametitle{Drop Redundant Columns}
Drop the columns we won't use:
\begin{itemize}
\item  The Age column since we will be using the AgeFill column instead.
\item     The SibSp and Parch columns since we will be using FamilySize instead.
\item     The PassengerId column since it won't be used as a feature.
\end{itemize}
\begin{lstlisting}
df_train = df_train.drop(['Age', 'SibSp', 'Parch', 'PassengerId'], axis=1)
\end{lstlisting}
\end{frame}

%%%%%%%%%%%%%%%%%%%%%%%%%%%%%%%%%%%%%%%%%%%%%%%%%%%
\begin{frame}[fragile]\frametitle{To NumPy Array}
Convert the DataFrame to a numpy array (columns: Survived, Pclass, Fare, Sex\_Val, AgeFill, FamilySize):
\begin{lstlisting}
>>> train_data = df_train.values
>>> train_data
array([[ 0.    ,  3.    ,  7.25  ,  1.    , 22.    ,  1.    ],
       [ 1.    ,  1.    , 71.2833,  0.    , 38.    ,  1.    ],
       [ 1.    ,  3.    ,  7.925 ,  0.    , 26.    ,  0.    ],
       ...,
       [ 0.    ,  3.    , 23.45  ,  0.    , 21.5   ,  3.    ],
       [ 1.    ,  1.    , 30.    ,  1.    , 26.    ,  0.    ],
       [ 0.    ,  3.    ,  7.75  ,  1.    , 32.    ,  0.    ]],
      shape=(891, 6))
\end{lstlisting}
\end{frame}

%%%%%%%%%%%%%%%%%%%%%%%%%%%%%%%%%%%%%%%%%%%%%%%%%%%%%%%%%%%%%%%%%%%%%%%%%%%%%%%%%%
\begin{frame}[fragile]\frametitle{}
\begin{center}
{\Large Modeling: Random Forest}
\end{center}
\end{frame}

%%%%%%%%%%%%%%%%%%%%%%%%%%%%%%%%%%%%%%%%%%%%%%%%%%%
\begin{frame}[fragile]\frametitle{Create the Model}
Create the random forest object (\texttt{random\_state=0} makes the results repeatable):
\begin{lstlisting}
from sklearn.ensemble import RandomForestClassifier

clf = RandomForestClassifier(n_estimators=100, random_state=0)
\end{lstlisting}
\end{frame}

%%%%%%%%%%%%%%%%%%%%%%%%%%%%%%%%%%%%%%%%%%%%%%%%%%%
\begin{frame}[fragile]\frametitle{Fit and Score}
Fit the training data and create the decision trees:
\begin{lstlisting}
# Training data features, skip the first column 'Survived'
train_features = train_data[:, 1:]

# 'Survived' column values
train_target = train_data[:, 0]

# Fit the model to our training data
clf = clf.fit(train_features, train_target)
score = clf.score(train_features, train_target)

# In the Python shell:
>>> "Mean accuracy of Random Forest: {0}".format(score)
'Mean accuracy of Random Forest: 0.9797979797979798'
\end{lstlisting}
\end{frame}

%%%%%%%%%%%%%%%%%%%%%%%%%%%%%%%%%%%%%%%%%%%%%%%%%%%
\begin{frame}[fragile]\frametitle{Fit and Score: Intuition}
\begin{block}{Intuition}
98\% looks great, but this is the model scored on the very same data it was trained on: a
random forest can nearly memorize its training set. It says little about how well the
model generalizes. The real test comes later, on data the model has never seen.
\end{block}
\end{frame}

%%%%%%%%%%%%%%%%%%%%%%%%%%%%%%%%%%%%%%%%%%%%%%%%%%%%%%%%%%%%%%%%%%%%%%%%%%%%%%%%%%
% Optional (Sep 2026 review): divider for the test-set predictions.
% \begin{frame}[fragile]\frametitle{}
% \begin{center}
% {\Large Predicting}
% \end{center}
% \end{frame}


%%%%%%%%%%%%%%%%%%%%%%%%%%%%%%%%%%%%%%%%%%%%%%%%%%%
% Optional (Sep 2026 review): the Kaggle submission is commented out, so predicting the Kaggle test set led nowhere; restore these three together with it.
% \begin{frame}[fragile]\frametitle{Random Forest: Load Test Data}
% Read the test data:

% \begin{center}
% \includegraphics[width=0.8\linewidth,keepaspectratio]{titk23}
% \end{center}

% Note the test data does not contain the column 'Survived', we'll use our trained model to predict these values.

% \begin{lstlisting}
% df_test = pd.read_csv("titanic_test.csv")
% df_test.head()
% \end{lstlisting}

% \end{frame}

%%%%%%%%%%%%%%%%%%%%%%%%%%%%%%%%%%%%%%%%%%%%%%%%%%%
% Optional (Sep 2026 review): the Kaggle submission is commented out, so predicting the Kaggle test set led nowhere; restore these three together with it.
% \begin{frame}[fragile]\frametitle{Random Forest: The clean\_data() Function}
% \texttt{clean\_data()} repeats every cleaning step, so the test set is treated exactly like the training set. The test set also has one missing Fare, filled here with the median.
% \begin{lstlisting}
% def clean_data(df, drop_passenger_id):
%     sexes = sorted(df['Sex'].unique())
%     df['Sex_Val'] = df['Sex'].map(dict(zip(sexes, range(len(sexes))))).astype(int)
%     df['AgeFill'] = df.groupby(['Sex_Val', 'Pclass'])['Age'] \
%                       .transform(lambda x: x.fillna(x.median()))
%     df['Fare'] = df['Fare'].fillna(df['Fare'].median())
%     df['FamilySize'] = df['SibSp'] + df['Parch']
%     df = df.drop(['Name', 'Sex', 'Ticket', 'Cabin', 'Embarked',
%                   'Age', 'SibSp', 'Parch'], axis=1)
%     if drop_passenger_id:
%         df = df.drop(['PassengerId'], axis=1)
%     return df
% \end{lstlisting}
% \end{frame}


%%%%%%%%%%%%%%%%%%%%%%%%%%%%%%%%%%%%%%%%%%%%%%%%%%%
% Optional (Sep 2026 review): the Kaggle submission is commented out, so predicting the Kaggle test set led nowhere; restore these three together with it.
% \begin{frame}[fragile]\frametitle{Random Forest: Clean Test Data}
% Clean the test data the same way as the training data:
% \begin{lstlisting}
% # Data wrangle the test set and convert it to a numpy array
% df_test = clean_data(df_test, drop_passenger_id=False)
% test_data = df_test.values

% # Get the test data features, skipping the first column 'PassengerId'
% test_x = test_data[:, 1:]

% # Predict the Survival values for the test data
% test_y = clf.predict(test_x)
% \end{lstlisting}
% \end{frame}

%%%%%%%%%%%%%%%%%%%%%%%%%%%%%%%%%%%%%%%%%%%%%%%%%%%%%%%%%%%%%%%%%%%%%%%%%%%%%%%%%%
% Commented out (Sep 2026): the Kaggle submission section is not needed for this session.
% \begin{frame}[fragile]\frametitle{}
% \begin{center}
% {\Large Kaggle Submission}
% \end{center}
% \end{frame}

%%%%%%%%%%%%%%%%%%%%%%%%%%%%%%%%%%%%%%%%%%%%%%%%%%%
% Part of the commented-out Kaggle submission section.
% \begin{frame}[fragile]\frametitle{Random Forest: Prepare for Kaggle Submission}
% Create a DataFrame by combining the index from the test data with the output of predictions, then write the results to the output:
% \begin{lstlisting}
% df_test['Survived'] = test_y
% df_test[['PassengerId', 'Survived']] \
%     .to_csv('../data/titanic/results-rf.csv', index=False)
% \end{lstlisting}
% \end{frame}

%%%%%%%%%%%%%%%%%%%%%%%%%%%%%%%%%%%%%%%%%%%%%%%%%%%
\begin{frame}[fragile]\frametitle{}
\begin{center}
{\Large Evaluating the Model}
\end{center}
\end{frame}

%%%%%%%%%%%%%%%%%%%%%%%%%%%%%%%%%%%%%%%%%%%%%%%%%%%
\begin{frame}[fragile]\frametitle{Train/Test Split}
Get an idea of accuracy on data the model has not seen.
We'll split our training data, 80\% will go to ``train'' and 20\% will go to ``test'':
\begin{lstlisting}
from sklearn import metrics
from sklearn.model_selection import train_test_split

train_x, test_x, train_y, test_y = train_test_split(train_features,
                                                    train_target,
                                                    test_size=0.20,
                                                    random_state=0)

# In the Python shell:
>>> print(train_features.shape, train_target.shape)
(891, 5) (891,)
>>> print(train_x.shape, train_y.shape)
(712, 5) (712,)
>>> print(test_x.shape, test_y.shape)
(179, 5) (179,)
\end{lstlisting}
\end{frame}

%%%%%%%%%%%%%%%%%%%%%%%%%%%%%%%%%%%%%%%%%%%%%%%%%%%
\begin{frame}[fragile]\frametitle{Held-Out Accuracy}
Use the new training data to fit the model, predict, and get the accuracy score:
\begin{lstlisting}
clf = clf.fit(train_x, train_y)
predict_y = clf.predict(test_x)

from sklearn.metrics import accuracy_score

# In the Python shell:
>>> print("Accuracy = %.2f" % accuracy_score(test_y, predict_y))
Accuracy = 0.84
\end{lstlisting}
\end{frame}

%%%%%%%%%%%%%%%%%%%%%%%%%%%%%%%%%%%%%%%%%%%%%%%%%%%
\begin{frame}[fragile]\frametitle{Quick Check: Why Did Accuracy Drop?}
\begin{block}{Think About It}
The random forest scored 98\% accuracy when evaluated on the full training data, but only
84\% on the held-out 20\% split. Why the big gap, and which number should you actually
trust to predict how the model will do on Kaggle's real test set?
\begin{enumerate}
  \item[A.] The 84\% number is wrong; the model is really 98\% accurate, the split just got
  unlucky.
  \item[B.] The 98\% reflects memorization of data the model already saw during training; the
  84\% reflects performance on unseen data, which is what actually matters.
  \item[C.] Random forests are inherently unstable, so both numbers are meaningless.
  \item[D.] The gap is caused by the confusion matrix, not by the train/test split.
\end{enumerate}
\end{block}
\end{frame}

%%%%%%%%%%%%%%%%%%%%%%%%%%%%%%%%%%%%%%%%%%%%%%%%%%%%%%%%%%%%%%%%%%%%%%%%%%%%%%%%%%
\begin{frame}[fragile]\frametitle{Quick Check: Why Did Accuracy Drop?}
\begin{block}{Answer}
\textbf{Correct: B.} A random forest with 100 trees has plenty of capacity to fit quirks of
the exact 891 training rows it saw, so scoring it on that same data mostly measures
memorization, not learning. The 84\% figure, measured on the 179 rows held out of training,
is the honest estimate of how the model performs on data it hasn't memorized: closer to what
Kaggle's actual (unseen) test set will show.
\end{block}
\end{frame}


%%%%%%%%%%%%%%%%%%%%%%%%%%%%%%%%%%%%%%%%%%%%%%%%%%%
\begin{frame}[fragile]\frametitle{Confusion Matrix}
View the Confusion Matrix: rows are the actual class, columns the predicted class.
\begin{itemize}
\item 101 and 49 on the diagonal: correct predictions
\item 9 died but were predicted to survive; 20 survived but were predicted to die
\end{itemize}
\begin{lstlisting}
confusion_matrix = metrics.confusion_matrix(test_y, predict_y)

print ("          Predicted")
print ("         |  0  |  1  |")
print ("         |-----|-----|")
print ("       0 | %3d | %3d |" % (confusion_matrix[0, 0],
                                   confusion_matrix[0, 1]))
print ("Actual   |-----|-----|")
print ("       1 | %3d | %3d |" % (confusion_matrix[1, 0],
                                   confusion_matrix[1, 1]))
print ("         |-----|-----|")

# printed:
#           Predicted
#          |  0  |  1  |
#          |-----|-----|
#        0 | 101 |   9 |
# Actual   |-----|-----|
#        1 |  20 |  49 |
#          |-----|-----|
\end{lstlisting}
\end{frame}


%%%%%%%%%%%%%%%%%%%%%%%%%%%%%%%%%%%%%%%%%%%%%%%%%%%
\begin{frame}[fragile]\frametitle{Classification Report}
Display the classification report:
\begin{itemize}
\item Precision: of the passengers predicted to survive, the share who did (49 of 58 = 0.84)
\item Recall: of the passengers who did survive, the share we found (49 of 69 = 0.71)
\item F1-score: one number that balances precision and recall; support: how many passengers of that class
\end{itemize}
\begin{lstlisting}
>>> from sklearn.metrics import classification_report
>>> print(classification_report(test_y, predict_y,
...                             target_names=['Not Survived', 'Survived']))
              precision    recall  f1-score   support

Not Survived       0.83      0.92      0.87       110
    Survived       0.84      0.71      0.77        69

    accuracy                           0.84       179
   macro avg       0.84      0.81      0.82       179
weighted avg       0.84      0.84      0.83       179
\end{lstlisting}
\end{frame}

%%%%%%%%%%%%%%%%%%%%%%%%%%%%%%%%%%%%%%%%%%%%%%%%%%%
\begin{frame}[fragile]\frametitle{Quick Check: Reading the Classification Report}
\begin{block}{Think About It}
Recall is 0.92 for ``Not Survived'' but only 0.71 for ``Survived'', even though both classes
get the same model. What's the most likely reason the model is worse at correctly finding
survivors specifically?
\begin{enumerate}
  \item[A.] The model is broken and should be retrained from scratch.
  \item[B.] ``Survived'' has fewer training examples (roughly 38\% of the 891 passengers), so
  the model has seen fewer patterns for that class and defaults toward the majority class
  more often.
  \item[C.] Recall and precision are unrelated to how much training data each class has.
  \item[D.] The 179-row test split happened to contain zero survivors.
\end{enumerate}
\end{block}
\end{frame}

%%%%%%%%%%%%%%%%%%%%%%%%%%%%%%%%%%%%%%%%%%%%%%%%%%%%%%%%%%%%%%%%%%%%%%%%%%%%%%%%%%
\begin{frame}[fragile]\frametitle{Quick Check: Reading the Classification Report}
\begin{block}{Answer}
\textbf{Correct: B.} This is a mild class-imbalance effect: ``support'' in the report shows
110 Not-Survived rows versus only 69 Survived rows in this test split, mirroring the full
dataset's roughly 62/38 split. With fewer survivor examples to learn from, the model is more
prone to guessing ``Not Survived'', which raises that class's recall while lowering
survivors' recall. This is the same class-imbalance idea worth watching for whenever one
outcome is rarer than the other.
\end{block}
\end{frame}

%%%%%%%%%%%%%%%%%%%%%%%%%%%%%%%%%%%%%%%%%%%%%%%%%%%
\begin{frame}[fragile]\frametitle{Cross-Validation}
\begin{itemize}
\item One 80/20 split gives one number, and it depends on which rows landed in the test set
\item Cross-validation repeats the split: train on 4 parts, test on the 5th, and rotate 5 times
\item Report the mean and the spread of the 5 scores
\end{itemize}
\begin{lstlisting}
from sklearn.model_selection import cross_val_score

clf = RandomForestClassifier(n_estimators=100, random_state=0)
scores = cross_val_score(clf, train_features, train_target, cv=5)

# In the Python shell:
>>> print(scores.round(3))
[0.777 0.803 0.854 0.803 0.848]
>>> print("Accuracy: %.3f +/- %.3f" % (scores.mean(), scores.std()))
Accuracy: 0.817 +/- 0.030
\end{lstlisting}
\end{frame}

%%%%%%%%%%%%%%%%%%%%%%%%%%%%%%%%%%%%%%%%%%%%%%%%%%%
\begin{frame}[fragile]\frametitle{Which Features Matter?}
\begin{itemize}
\item A fitted random forest reports how much each feature helped its splits
\item Here Fare, Sex\_Val and AgeFill carry most of the weight; Pclass and FamilySize carry little
\item Caution: these importances favor continuous features, so treat them as a rough guide
\end{itemize}
\begin{lstlisting}
clf = clf.fit(train_features, train_target)
names = ['Pclass', 'Fare', 'Sex_Val', 'AgeFill', 'FamilySize']

# In the Python shell:
>>> for name, value in zip(names, clf.feature_importances_):
...     print("%-10s %.3f" % (name, value))
...
Pclass     0.083
Fare       0.294
Sex_Val    0.275
AgeFill    0.267
FamilySize 0.081
\end{lstlisting}
\end{frame}

%%%%%%%%%%%%%%%%%%%%%%%%%%%%%%%%%%%%%%%%%%%%%%%%%%%%%%%%%%%%%%%%%%%%%%%%%%%%%%%%%%
\begin{frame}[fragile]\frametitle{}
\begin{center}
{\Large Next?}
\end{center}
\end{frame}

%%%%%%%%%%%%%%%%%%%%%%%%%%%%%%%%%%%%%%%%%%%%%%%%%%%
\begin{frame}[fragile]\frametitle{To Improve Accuracy: Feature Work}
Feature Preparation, Selection, and Engineering
\begin{itemize}
\item How to determine which features in your model are the most-relevant to your predictions
\item Ways to reduce the number of features used to train your model and avoid over-fitting
\item Techniques to create new features to improve the accuracy of your model
\end{itemize}
\end{frame}

%%%%%%%%%%%%%%%%%%%%%%%%%%%%%%%%%%%%%%%%%%%%%%%%%%%
\begin{frame}[fragile]\frametitle{To Improve Accuracy: Model Tuning}
Model Selection and Tuning
\begin{itemize}
\item 
    How other classification algorithms work
\item About hyper-parameters, and how to select the hyper-parameters that give the best predictions
\item  How to compare different algorithms to improve the accuracy of your predictions
\end{itemize}
\end{frame}

